% 01 · Introduction: the science question and the plan of the workbook
\section{INTRODUCTION}
\label{sec:intro}

\deckslides{slides 1--3 (the science question, the data, and the clustering
problem) and slide 51 for the reference list.}

Every star that exists today was born in a cloud of gas and dust, and most of
them were born in the company of thousands of siblings. A cluster forms from
one well-mixed reservoir, so its members inherit almost the same chemical
composition: element ratios that agree to better than 0.1 dex, far below the
spread between one cluster and the next. Then the cluster dissolves. Tidal
fields, encounters with giant molecular clouds and the shear of the Galactic
disc peel the members away one by one, and after a few hundred million years
(billions, for the most massive clusters) the stars are scattered across
kiloparsecs of the disc, moving on orbits that no longer have anything in
common. Their chemistry, however, is frozen. A star keeps the composition it
was born with, because nothing in a stellar interior efficiently mixes the
surface layers. \emph{Chemical tagging} is the attempt to use that frozen
fingerprint to put dispersed siblings back together: to look at a field of
stars, all of them currently unrelated, and work out which of them were born in
the same lost cluster \citep{Freeman:02, BlandHawthorn:16}.

This workbook is the long-form companion to the IAA-SO 2026 lecture on that
problem. The lecture is a tour of the unsupervised-learning toolbox with one
astronomical question attached to it; the workbook goes deeper on both halves:
the mathematics of each algorithm, and the measurements we made with it on 25
real open and globular clusters.

\begin{marginnote}[]
\entry{Chemical tagging}{Using the element abundances of stars to recover
which of them were born together, long after the birth cluster has dispersed.}
\entry{C-space}{The abundance space in which tagging happens: here 16
elements measured from a spectrum, one point per star.}
\end{marginnote}

\subsection{A cluster dissolves; its chemistry does not}

The premise is a physical fact about the interstellar medium: mixing is fast
and efficient, so a cluster of $10^3$--$10^5$ stars forms from gas that was
chemically homogeneous to a fraction of a dex. Those stars also formed in one
burst, so they share an age. What they do \emph{not} share, once the cluster
is gone, is position or velocity. Galactic archaeology therefore faces a
reconstruction problem: the cluster is a fossil, and chemistry is the only
part of it that is still legible.

Two things make the reconstruction hard, and both are properties of the data
rather than of the algorithms.

\begin{enumerate}
\item \textbf{Precision.} The abundance differences between clusters of similar
age and birth radius are small. At $0.05$--$0.1$ dex per element (roughly the
precision of a modern survey pipeline) the chemical differences between two
coeval solar-metallicity open clusters can be smaller than the measurement
noise for the elements in question \citep{GarciaDias:19, Casamiquela:21}.
Measured differentially (same instrument, same lines, cluster members only) the
internal coherence improves to $0.03$ dex, and the cluster-to-cluster
differences do not improve with it: the gas the clusters formed from was
already mixed to a comparable level, $0.02$--$0.03$ dex of scatter in nearby
spiral discs, correlated over scales below $600$\,pc \citep{Kreckel:20}. The
same efficient mixing that makes one cluster chemically coherent is what limits
how different two clusters can be.
\item \textbf{Contamination.} The field is crowded. In the sample used here,
  1\,002 stars belong to the 25 target clusters and 357\,056 do not: roughly
  one member in every 300 stars. Any search that returns a group has to be
  judged on both how many real members it found (recall) and how many of the
  stars in the group are real members (precision).
\end{enumerate}

\subsection{The task is clustering, not classification}

Chemical tagging is a \emph{clustering} problem: we are not handing labels to
a supervised model, we are asking an unsupervised algorithm to tell us how
many groups the data contains and which star belongs to which. That
distinction has consequences which reappear throughout the workbook.

\begin{issues}[WHY IT MATTERS THAT NOTHING IS LABELLED]
\begin{enumerate}
\item There is no ground truth inside the algorithm. A supervised classifier
  is scored on held-out labels; a clustering algorithm cannot be, so a
  clustering result is only as good as the validation you bring to it from
  outside (\textbf{Section~\ref{sec:validation}}).
\item Hyperparameters must be chosen without labels, which is why the method
  chapters spend as much time on \emph{what the knob means} as on the
  mathematics that uses it.
\item Feature choice is part of the model. In this problem the same algorithm
run on abundances, on a spectral latent, or on kinematics answers different
questions, and concatenating two such spaces does not combine their strengths,
it dilutes the smaller one (\textbf{Section~\ref{sec:benchmark}}).
\end{enumerate}
\end{issues}

\subsection{Two published verdicts, and the one this workbook adds}

The literature is not unanimous, and the disagreement is instructive because it
is not about methods. It is about tasks and precision.

\citet{Kos:17} ran t-SNE on GALAH abundances around the Pleiades and reported
success: seven of nine clusters recovered, plus two new Pleiades members six
degrees from the cluster centre. Their test case was a young cluster
($\sim\!100$ Myr), their verification was kinematic, and their grouping step
was visual inspection of the embedding.

\citet{GarciaDias:19} asked the same question with APOGEE abundances over
23 clusters and returned a sceptical answer: some clusters are chemically
indistinguishable from each other, and abundances alone are enough to
identify \emph{families} of clusters rather than individual ones. The
reconciliation is not that one study is wrong. The two tagged different
objects, at different ages, with different abundance precision, and asked
different things of the answer.

A third result belongs on the table too, because it cuts against the
pessimistic reading. \citet{Hogg:16} clustered 15 \emph{Cannon} abundances for
$\sim\!10^5$ APOGEE stars with K-means, using no positional information at all,
and showed that abundance-space overdensities really are phase-space clusters.
Their abundances reach $\sim\!0.04$ dex, comfortably inside the scatter of the
pipeline used here. The scope of any negative result in this field is therefore
\emph{these abundances, these clusters}, never ``abundances cannot tag''.

The current state of the art sits in between. \citet{Spina:25} trained a graph
attention autoencoder on chemistry, orbital similarity and age for
$\sim\!47\,000$ APOGEE thin-disc stars, obtained 282 groups and recovered five
of six open clusters. It is the natural competitor to this work and we have
not benchmarked against it; \textbf{Section~\ref{sec:spectral}} says what we
expect a like-for-like comparison to change.

\begin{marginnote}[] \entry{Strong tagging}{Recovering a specific lost birth
cluster from chemistry alone: the hard version, and the one this workbook
benchmarks.} \entry{Weak tagging}{Grouping stars into chemical families.
Populations that share a formation epoch rather than a common birth site. The
realistic outcome at current precision.}
\end{marginnote}

\subsection{What this workbook measures}

The second half is a benchmark we ran ourselves. We take 25 clusters (18 open,
7 globular) with kinematic membership from Gaia DR3, cut the APOGEE DR19 sample
around them, and ask three families of method (t-SNE, UMAP and EVoC) to recover
members from 16 abundances, from a 256-dimensional latent of the full spectrum,
or from kinematics, all scored against the same kinematic truth. On top of that
we train a masked spectral autoencoder that never sees an abundance label and
ask whether it tags better than the abundances themselves.

Four findings organise the results, each developed in its own section:

\begin{enumerate}
\item \textbf{Abundances narrow; kinematics decide.} Abundance-space
  clustering separates the 25 clusters only weakly (homogeneity $0.56$ for the
  best projector) where kinematics reach $0.94$, and adding abundances to
  kinematics never helps.
\item \textbf{Full-sky collapses; regions work.} Embedding $163\,000$ stars at
  once buries every cluster in field stars. The published t-SNE tagging
  operated on $30$--$45^\circ$ regions, and reproduces only there.
\item \textbf{The full spectrum beats the element ratios.} A self-supervised
latent of 8\,575 spectral pixels separates the same stars better than the 16
ASPCAP abundances, and the gap survives removing the globular cluster that
dominates the sample. A linear compression (PCA) sits in between, which is the
honest way to state the result.
\item \textbf{Chemical tagging is a two-survey problem.} APOGEE sees the
  giants that pin the distance; GALAH sees the main sequence that pins the
  age. Neither survey alone recovers both cluster parameters.
\end{enumerate}

\subsection{How to read this workbook}

Deck and workbook are complements. The deck carries the animations (K-means
converging, DBSCAN growing, the density sweep, the t-SNE descent, the UMAP
layout, the PLSCAN persistence barcode) and the workbook carries the equations,
the hyperparameter reasoning and the measured numbers. Where a deck figure
moves, the workbook prints a filmstrip of key frames and says so in the
caption.

\begin{figure}[htbp]\centering
  \begin{tabular}{@{}c@{\hskip18pt}c@{}}
    \fig[0.28]{qr_deck.png} &
    \fig[0.28]{qr_repo.png}
  \end{tabular}
  \caption{Left: the live deck, where the animated versions of the figures in
  this workbook run. Right: the repository that produced every number quoted
  in \textbf{Sections~\ref{sec:benchmark}} and \ref{sec:spectral}. Both are
  unlisted but openly reachable.}
  \label{fig:qrcodes}
\end{figure}

The chapters follow the order of the lecture. \S\ref{sec:data} is the data;
\S\ref{sec:fundamentals}--\S\ref{sec:evoc} are the toolbox, each method
motivated by the failure of the previous one: K-means
(\S\ref{sec:kmeans}), the $k$-nearest-neighbour primitive
(\S\ref{sec:knn}), DBSCAN (\S\ref{sec:dbscan}),
HDBSCAN$^*$ (\S\ref{sec:hdbscan}), PLSCAN (\S\ref{sec:plscan}), validation
(\S\ref{sec:validation}), t-SNE (\S\ref{sec:tsne}), UMAP
(\S\ref{sec:umap}) and EVoC (\S\ref{sec:evoc}).
\S\ref{sec:benchmark}--\S\ref{sec:physics} are the measurements, and
\S\ref{sec:assignment} is the assignment.

\begin{issues}[EXERCISES]
\begin{enumerate}
\item A cluster of $10^4$ stars dissolves over a few hundred Myr; a globular
  cluster of $10^6$ stars survives for more than 10 Gyr. Explain from that why
  chemical tagging should be \emph{easier} for old metal-poor clusters than
  for young solar-metallicity ones, then check your answer against the
  per-cluster results in \textbf{Section~\ref{sec:benchmark}}.
\item The field-to-member ratio here is about 300:1. If a search returns 100
  candidates for a cluster with 20 true members, what are its precision and
  recall? Which would you rather sacrifice, and what does that choice cost
  downstream?
\item Read \citet{Hogg:16} and \citet{Casamiquela:21} side by side: one is
  optimistic about strong chemical tagging and one is not. Write two sentences
  each on which \emph{measurement} they disagree about.
\end{enumerate}
\end{issues}
